\documentclass[10pt,a4paper]{amsart}
\usepackage[margin=19mm]{geometry}
\usepackage{amsmath,amssymb,mathtools}
\usepackage[scr=rsfs,cal=euler]{mathalfa}
\usepackage{xfrac}
\usepackage[unicode,hidelinks,bookmarks=false]{hyperref}
\newcommand{\ie}{i.e.\ }
\newcommand{\cf}{cf.\ }
\newcommand{\resp}{respectively\ }
\newcommand{\Subsection}{\subsection}
\providecommand{\nobreakdash}{\nobreak\hbox{-}}
\setlength{\emergencystretch}{2em}
\multlinegap0pt
\usepackage{mdframed}
\usepackage{mdframed}
\usepackage{mdframed}
\usepackage{mdframed}
\usepackage{amssymb}
\usepackage{xcolor}
\usepackage{cancel}
\usepackage{mdframed}
\usepackage{mathtools}
\usepackage{enumerate}
\usepackage[shortlabels]{enumitem}
\usepackage{tikz}
\usepackage{bm}
\newtheorem{lemma}{Lemma}
\newtheorem{theorem}{Theorem}
\renewcommand{\bmod}[1]{\ ( \mathrm{mod} \, #1 )}
\title{On the digits of the sum of proper divisors}
\author{K\"ubra Benl\.i, C\'ecile Dartyge, Charlotte Dombrowsky, Paul Pollack,, Lola Thompson}
\date{Source: arXiv:2607.18981v1 (CC BY 4.0)}
\begin{document}
\maketitle

\noindent\emph{Source and licence.} On the digits of the sum of proper divisors. Version arXiv:2607.18981v1 (CC BY 4.0), distributed by arXiv under the Creative Commons Attribution 4.0 International licence, http://creativecommons.org/licenses/by/4.0/, as recorded in the arXiv licence metadata for this exact version and shown on the abstract page https://arxiv.org/abs/2607.18981v1. Full licence notice: \texttt{licenses/arxiv-2607.18981-CC-BY-4.0.txt}.\par\smallskip

\noindent\textit{Editorial scope.} Counted unit: the article's theorem thm:firstk (source0.tex lines 364--367). The article's complete original proof of it is delivered with it (source0.tex lines 368--411, one \texttt{proof} environment of 44 lines). No mathematics is written, repaired or re-derived: the statement and the proof are the article's own lines, in the article's own notation, and no retained line is substituted or re-flowed.  Retained uncounted as the source-local context the proof needs: lines 300--305.  Nothing was omitted from any retained span, so the package carries no omission record.  No retained line contains a citation, so the wrapper carries no bibliography and no bibliography entry.  No retained line contains a figure environment, an image-inclusion command or drawing code, so no image is delivered and the package declares no asset dependency.  Editorial formatting only, in addition: an AMS \texttt{amsart} preamble that restates the article's own declarations and macros used by the retained lines, namely the environments \texttt{lemma}, \texttt{theorem} on separate counters, together with the standard packages those lines need.  The retained environments print in this document as Lemma 1, Theorem 1 in order of appearance, whereas the article prints the same items with the numbering of its own full document.  The complete native source of the article as distributed in the arXiv e-print for this exact version is bundled unchanged as \texttt{sources/205695/source0.tex}.
\begin{lemma}\label{lem: n to nz}
	Let $x$ be large and $k\in \mathbb{N}$ with $k(x)\to\infty$ as $x\to \infty$. Put $z = g^{10k}$.
	For all but $o(x)$ values of $n \leq x$, one has 
	$0 \leq s(n)/n - s(n_z)/n_z < z^{-1/2}$
	and $s(n_z)/n_z \geq 1/k.$
	\end{lemma}

	\begin{theorem}\label{thm:firstk}
	Assume $k(x)\to \infty$ as $x\to \infty$. 	
	Then the number of $n\leq x$ where, in the base $g$ expansion of $s(n)$, not all digits are represented among the first $k$ digits is $o(x)$. In other words, asymptotically $100\%$ of integers $n\leq x$ are such that $s(n)$ has all $g$ digits in base $g$ appearing among its first (i.e., leftmost) $k(x)$ digits.
	\end{theorem}

	\begin{proof}
    We can assume our function $k=k(x) \to \infty$ grows very slowly, e.g., $k(x) \leq \log\log\log x$ (otherwise, we can simply replace $k(x)$ with $\min\{k(x),\log\log\log x\}$).
     
Let $n\leq x$ be such that in the base $g$ expansion of $s(n)$, not all digits are represented among the first $k$ digits. Let $r\in \mathbb{N}$ be such that \[ g^{r-1}\leq n<g^r. \] By our first remark above, we may assume that \[ g^r>\frac{x}{g^{\log\log k}}.\] Let $m\in \mathbb{N}$ be such that  \[ g^{m-1}\leq s(n)<g^m.\] By our second remark above, we may assume  that 
\[ |m-r| < \log \log k. \] 

Write $s(n) = g^{m-k} L  + T$, where $L, T\in \mathbb{N}$ with $0\leq T< g^{m-k}$. Then $L < g^{k}$, and the digits of $L$ are the $k$ leading digits of $s(n)$. Note that the number of possibilities for $L$ is $\ll g^{ck}$, for a constant $c < 1$ depending only on $g$ (because $L$ is missing at least one base $g$ digit). This puts $s(n)$ in the interval  \begin{align*} I_L&:=[L\cdot g^{m-k},(L+1)\cdot g^{m-k}) \\&:=[t_L,T_L). \end{align*} 


We fix $r$, $m$, and $L$, and count the number of corresponding $n$. It is convenient to sort these $n$ into residue classes modulo $M = \prod_{p \leq z} p^{z}$, where $z = g^{10k}$. By the Prime Number Theorem, $M \leq 3^{z^2} =x^{o(1)}$, where the last equality comes from our restriction on the size of $k$. Let  $n'\in (0,M]$ be the least positive representative of $n \bmod M$  and $n_{z}$ be the $z$-smooth part of $n'$. Excluding the $o(x)$ exceptional values of $n$ described in Lemma \ref{lem: n to nz} above, we have $$\frac{s(n)}{n} \leq \frac{s(n_z)}{n_z} \left(1 +\frac{n_z}{s(n_z)}  z^{-1/2}\right) \leq {s(n_z)} (1 + k z^{-1/2}).$$ 
    
We now fix also $n'\in \mathbb{N}\cap (0,M]$ and consider only nonexceptional integers $n \equiv n' \bmod M$ corresponding to our  choices of $r$,$m$, and $L$ from above.

If $s(n) = n \cdot s(n)/n$ is in one of the intervals $I_L$ from above, then $n$ is in $(n/s(n)) I_L$. Notice that $s(n)/n \geq s(n_z)/n_z$, and so
	$n/s(n)\leq  n_z/s(n_z)$. Since $T_L$ is the upper endpoint of $I_L$, the dilated interval $(n/s(n))I_L$ has upper endpoint at most \[ T_{L'} := (n_z/s(n_z)) T_L.\]
	
	Now, we consider the lower endpoint. We have
    
	$$\frac{s(n)}{n} \leq \frac{s(n_z)}{n_z}+ z^{-1/2}
	= \frac{s(n_z)}{n_z}\left(1 + z^{-1/2} \frac{n_z}{s(n_z)}\right)
\leq \frac{s(n_z)}{n_z} (1 + z^{-1/3})$$
since $z=g^{10k}.$ Hence,
 \[ \frac{n}{s(n)} \geq \frac{n_z}{s(n_z)} (1-z^{-1/3}).\] Thus, if $t_L$ is the lower endpoint of $I_L$, then $(n/s(n))I_L$ has lower endpoint at least \[ t_{L'} := n_z/s(n_z) (1-z^{-1/3}) t_L.\] 
 
Every nonexceptional $n\equiv n'\bmod{M}$ corresponding to our  choices of $r$,$m$,  and $L$ belongs to the interval $[t_{L'}, T_{L'})$. The length of this interval is
$$T_{L'}-t_{L'}=\frac{n_z}{s(n_z)} T_L-\frac{n_z}{s(n_z)}(1-z^{-1/3}) t_L=\frac{n_z}{s(n_z)} |I_L| + z^{-1/3} \frac{n_z}{s(n_z)} t_L,$$
and so
\begin{align*}
\#\{n\leq x: n \equiv n' \bmod M,\,\, t_{L'}\leq n< T_{L'}    \}&\leq \frac{1}{M}\frac{n_z}{s(n_z)} \left(|I_L| + z^{-1/3} t_L\right)+1
\\&\leq 1 + \frac{k g^{m-k}(1+z^{-1/3}L)}{M}.
\end{align*}

We bound the total number of nonexceptional $n$ by summing over parameters previously held fixed. First, we sum over the $M$ possible values of $n' \in (0,M]$. 
Recalling that $z=g^{10k}$ and $L\ll g^{ck}$ for $c<1$, the number of $n$ arising, for a fixed choice of $r, m$, and $L$, is at most
    \[ M + 2k \cdot g^{m-k} \le 3k g^{m-k}. \]
For the last inequality, we used that $M=x^{o(1)}$ while (for large $x$)
\[ g^{m-k} = g^{r} g^{m-r} g^{-k} \gg \frac{x}{g^{\log\log{k}}} \cdot \frac{1}{g^{\log\log{k}}} \cdot \frac{1}{g^k} > \frac{x}{g^{2k}} \ge \frac{x}{(\log\log{x})^{O(1)}}.\]
    
    
   
	
Next, we sum on the possibilities for $L$; this gives an upper bound of the form $3k \cdot g^{-c' k} g^m$, for a positive $c'>0$, bounded away from $0$. Furthermore, $m \leq r + \log \log k$. Summing $3k \cdot g^{-c' k} g^m$ on these $m$, we obtain a quantity bounded above by $g^r g^{-c'' k}$, for another small $c'' > 0$. Summing on the possibilities for $r$ gives an upper bound of $x g^{-c''' k}$ for $c'''>0$. If we combine this upper bound with the previously mentioned $o(x)$ contributions from the exceptions, we still obtain $o(x)$ for our final count.
		
	\end{proof}

\end{document}
